\section{Evaluation}
\label{sec:evaluation}

\subsection{Methodology}

We ask whether the coupled design improves end-to-end low-memory MoE decoding, whether
the gain transfers across GPU generations and memory budgets, and whether the measured
expert traffic supports the proposed mechanism. Table~\ref{tab:setup-fairness}
summarizes the controlled comparison.

\begin{table*}[t]
  \centering
  \caption{Evaluation matrix and controlled baseline contract. All methods use the
  same target checkpoints, prompts, output policy, runtime, and memory ceiling.}
  \label{tab:setup-fairness}
  \footnotesize
  \begin{minipage}[t]{0.51\textwidth}
  \textbf{(a) Target models}\\[-1pt]
  \begin{tabular*}{\linewidth}{@{\extracolsep{\fill}}lrrrr@{}}
    \toprule
    Target & Size (GiB) & Layers & Experts/layer & Top-$k$ \\
    \midrule
    Mixtral-8x7B & 87.0 & 32 & 8 & 2 \\
    DeepSeek-V2-Lite & 29.3 & 27 & 64 & 6 \\
    Qwen1.5-MoE & 26.7 & 24 & 60 & 4 \\
    OLMoE-1B-7B & 12.9 & 16 & 64 & 8 \\
    \bottomrule
  \end{tabular*}
  \end{minipage}%
  \hfill
  \begin{minipage}[t]{0.47\textwidth}
  \textbf{(b) Compared execution paths}\\[-1pt]
  \begin{tabular*}{\linewidth}{@{\extracolsep{\fill}}lll@{}}
    \toprule
    Method & Draft path & Verification supply \\
    \midrule
    Classic SubSpec & target W4 & full BF16 page-in \\
    SS-MoE & reduced-$k$ & native BF16 \\
    SpecMoE & pinned+remap & native BF16 \\
    Ours & shared target W4 & exact recovery \\
    \bottomrule
  \end{tabular*}
  \end{minipage}
\end{table*}

The headline uses an NVIDIA A100 and an 8~GiB software-enforced budget; cross-platform
experiments repeat the matched Qwen and DeepSeek matrix on RTX~4090 and RTX~5090,
while the deployment sweep uses 8, 12, 16, 20, and 24~GiB on RTX~4090. Four workloads cover
reasoning, code, long-input summarization, and dialogue: GSM8K
\cite{cobbe2021training}, HumanEval \cite{chen2021codex}, BillSum
\cite{kornilova2019billsum}, and MT-Bench first turn \cite{zheng2023mtbench}.

Each cell uses four fixed prompts and three fresh-process replicas with 128 timed
output tokens. Throughput counts committed, target-authorized tokens over the complete
prefill--draft--verify--commit path. We first pair every method with its Classic
MoE-SubSpec run on the same prompt and replica, then report workload and model
geometric means; error bars use a paired hierarchical 95\% bootstrap. The headline
contains 180 complete execution bundles and 528 matched prompt ratios.

Correctness is established before timing: registered corrections must reconstruct
the source BF16 bit pattern, recovered experts must match the native expert operator,
and verification must overwrite provisional KV state while preserving the clean
prefix. The target router, native BF16 operator, and greedy acceptance rule are
unchanged; any undeclared full-BF16 fallback invalidates an exact run.

\begin{figure*}[t]
  \centering
  \includegraphics[width=\textwidth]{f2_classic_normalized_headline.pdf}
  \caption{End-to-end throughput at 8~GiB, normalized to matched Classic
  MoE-SubSpec. Bars group four workloads within each model; dashed lines mark the
  Classic baseline, error bars show paired uncertainty, and labels report per-model
  geometric means. Mixtral--SpecMoE is OOM under the common budget.}
  \label{fig:f2}
\end{figure*}

\subsection{End-to-end performance}

Figure~\ref{fig:f2} shows that the coupled system consistently converts compact draft
state into target-side speedup. Across all four model families, \system achieves a
1.96$\times$ geometric mean over Classic MoE-SubSpec. The gains are
2.57$\times$ on Mixtral, 2.12$\times$ on DeepSeek, 1.64$\times$ on Qwen, and
1.65$\times$ on OLMoE. The result spans reasoning, code, summarization, and dialogue
rather than depending on a single prompt style.

The comparison also separates the benefit of speculation from the benefit of exact
cross-phase reuse. SS-MoE and SpecMoE accelerate selected cells, but \system attains
the highest model-level geometric mean in every family. On Mixtral, ordinary pinned
BF16 speculation cannot enter the shared 8~GiB envelope, whereas the W4 base and
compact correction representation retain an executable verification path. The
largest gains appear where the verifier would otherwise repeatedly move the widest
expert state, matching the system's design objective.

\textbf{Consistency across workloads.} Every completed \system bar in
Fig.~\ref{fig:f2} remains above its matched Classic run. This matters because the four
workloads exercise different sequence shapes: GSM8K emphasizes iterative reasoning,
HumanEval short code continuation, BillSum long-input summarization, and MT-Bench
dialogue. The gain is therefore not produced by one favorable prompt length. Nor is
it a claim that model size alone determines speedup. The relevant quantity is the
joint pressure of the routed union, expert footprint, and reusable W4 residency. The
per-model geometric means summarize that interaction without hiding the workload
distribution behind a single aggregate.

\textbf{A demanding baseline contract.} Classic MoE-SubSpec already uses the target's
quantized experts as a highly aligned training-free draft and verifies proposals with
the original target. It therefore captures the speculative benefit available before
our exact-reuse mechanism is enabled. Normalizing each prompt and replica to this
path asks the narrow systems question of the paper: once proposal quality and target
semantics are fixed, can one compact target state eliminate duplicated expert and
sequence movement? The consistent cell-level gains answer that question directly:
exact cross-phase reuse reduces duplicated state movement while preserving the same
proposal and verification semantics.

\begin{figure*}[t]
  \centering
  \includegraphics[width=0.98\textwidth]{f3_cross_platform_mechanism.pdf}
  \caption{Generality and mechanism evidence. (a) Workload and geometric-mean
  speedups over matched Classic MoE-SubSpec on three GPUs at 8~GiB. (b) The
  deployment envelope on RTX~4090 over 8, 12, 16, 20, and 24~GiB; shading spans workloads.
  (c) Increasing proposal width expands the verification expert union while draft
  coverage rises. Marker size reflects the full-BF16 traffic lower bound. (d) Exact
  recovery replaces full-BF16 movement with compact correction traffic and reuses
  resident W4 bases.}
  \label{fig:f3}
\end{figure*}

\subsection{Generality and mechanism}

\textbf{Across hardware.} Figure~\ref{fig:f3}(a) preserves the same mechanism across
three GPU generations. DeepSeek reaches 2.12$\times$, 1.25$\times$, and 1.48$\times$
on A100, RTX~5090, and RTX~4090; Qwen reaches 1.64$\times$, 1.82$\times$, and
1.31$\times$. Every platform--model geometric mean remains above Classic despite
different compute-to-bandwidth balances, showing that the benefit follows the state
hierarchy rather than one device-specific kernel profile.

\textbf{Across memory budgets.} Figure~\ref{fig:f3}(b) shows a broad deployment
envelope rather than a single capacity point. Across 8, 12, 16, 20, and 24~GiB,
DeepSeek delivers 1.48$\times$, 1.88$\times$, 1.77$\times$, 1.82$\times$, and
1.70$\times$; Qwen delivers 1.31$\times$, 1.72$\times$, 1.59$\times$,
1.63$\times$, and 1.49$\times$. The dense sweep exposes a stable benefit throughout
the tested interval and discrete residency transitions as more W4 bases and recovered
BF16 experts fit in HBM.

\textbf{Why it works.} Figure~\ref{fig:f3}(c) connects proposal width to the physical
MoE working set. As $\gamma$ increases through 1, 2, 4, and 8, draft coverage rises,
but the verification union also expands; a design that supplies every new member as a
complete BF16 expert gives back part of the speculative gain. Figure~\ref{fig:f3}(d)
shows the response: correction traffic is 38\% smaller than the corresponding full
BF16 movement, and 89\% of DeepSeek and 76\% of Qwen recovery bases are already
resident. One compact expert state therefore contributes twice---first to aligned
drafting, then to exact verification---while shared prefill and KV management remove
the remaining cross-phase duplication.

Taken together, the figures establish one causal story: speculative width creates a
larger expert-supply opportunity; target-derived W4 state makes that opportunity
reusable; exact correction preserves native verification; and unified residency plus
shared KV turns the representation advantage into end-to-end throughput across
models, workloads, budgets, and hardware platforms.

\textbf{Cross-figure triangulation.} The four panels of Fig.~\ref{fig:f3} measure
different links of that story rather than four unrelated optimizations. Panel (c)
shows the demand that verification creates; panel (d) shows how much of that demand
can be served from a smaller transferred representation and an already resident
base. Panel (b) then varies the capacity available to retain those views, and panel
(a) varies the machine balance that prices transfer and reconstruction. Their common
outcome---speedup above Classic across every platform--model mean---is the expected
signature of state reuse whose value changes in magnitude, but not in direction,
with the deployment environment.

\section{Discussion}
\label{sec:discussion}

\textbf{The contribution is the coupling, not an isolated fast path.} A quantized
draft alone still leaves the verifier to move complete experts, while compact expert
repair alone still pays for an independent draft and duplicated sequence state.
\system joins these choices around one target-derived object: the W4 expert generates
the proposal, survives as reusable residency, and becomes the base for exact BF16
materialization. The shared router and dense path align the two execution views;
unified W4/BF16 residency turns alignment into memory reuse; shared prefill and KV
management carry that reuse across the entire decoding round. Each mechanism removes
a different form of duplication, and their composition is what converts a smaller
representation into end-to-end throughput.

\textbf{Exactness is a systems enabler.} The verifier never treats the quantized
expert as authoritative. Instead, the correction representation restores the native
BF16 bit pattern before the existing expert operator runs. This clean boundary lets
the system reuse approximate state aggressively without changing the target router,
acceptance rule, or committed KV semantics. It also simplifies deployment: the
correction is generated from the same checkpoint as the W4 base, requires no
retraining, and can be verified independently of the speculative policy. Draft
quality may affect how much work is saved, but it does not alter which model defines
the output.

\textbf{The results expose a reusable deployment principle.} The same mechanism wins
on models with different expert counts and routing widths, on compute-heavy and
bandwidth-heavy GPUs, and across a substantial memory interval. The absolute balance
between drafting, reconstruction, and transfer changes across these settings, yet the
controller presents the same three resources---compact bases, materialized exact
experts, and one KV trajectory. A deployment can therefore tune residency capacities
to its device without introducing a new draft model or changing target semantics.
This is especially attractive for a serving stack that must support several MoE
families under heterogeneous memory limits.

The memory sweep also clarifies why the benefit need not vary monotonically with
capacity. Each budget changes which W4 bases and materialized BF16 views survive a
round boundary, so the effective transfer path changes at discrete residency
thresholds rather than continuously. What remains invariant is the ordering of the
two supply strategies: across the complete sweep, reusing a target-derived base and
moving only its correction retains an advantage over replacing that base with a full
expert. This behavior links the capacity study back to the representation choice,
rather than treating the sweep as an independent tuning result.

\textbf{The architecture is complementary to lower-level optimizations.} Expert
prefetch, asynchronous I/O, fused kernels, and topology-aware placement can accelerate
either the W4 admission or exact materialization path; proposal schedulers can improve
accepted length without changing the recovery contract. \system supplies the missing
state abstraction connecting these techniques: a compact expert has value beyond
drafting because it remains an exactification base, and a materialized BF16 expert is
a cacheable view rather than a second checkpoint. This separation of state identity
from execution precision opens a broader design space for memory-aware speculative
MoE serving.

\textbf{The resulting abstraction is operational, not merely conceptual.} A serving
runtime can expose each logical expert as one identity with two executable views:
compact W4 and materialized BF16. Admission decisions then operate on future value
rather than tensor size alone: retaining a W4 page preserves both draft capacity and
the ability to recreate an exact expert, while retaining BF16 removes immediate
recovery work. This is the key distinction from treating quantization as a separate
model-compression stage. The representation participates in scheduling, supply, and
correctness as one managed object.

\textbf{The evidence supports the full-system claim.} End-to-end throughput supplies
the outcome; the hardware and memory sweeps show portability; union growth identifies
the pressure point; transfer and residency measurements identify the saved resource;
and bit-exact reconstruction plus target-authorized commit preserve semantics. No
single measurement is asked to prove the entire mechanism. Together they show how a
target-derived draft, exact expert recovery, unified residency, and shared sequence
state reinforce one another. That composition---rather than an isolated compression
ratio or speculative policy---is the central systems contribution.

\FloatBarrier

\subsection{Serving implications}

The coupled representation changes what a memory budget buys. In a conventional
substitute path, draft capacity and exact-verification capacity are separate claims
on memory: evicting a draft expert destroys proposal capacity, while loading its
BF16 counterpart consumes a second, larger allocation. In \system, a retained W4
page preserves two options. It can execute immediately in the draft view, or combine
with a correction to recreate the exact view. The controller can therefore allocate
HBM among W4 bases, hot BF16 views, and shared KV according to observed route reuse,
without changing the checkpoint or speculative algorithm.

Finally, the state hierarchy provides a stable integration boundary. Existing MoE
kernels continue to consume BF16 experts; existing speculative logic continues to
consume target logits and commit a target-approved prefix. The new mechanism sits
between these layers as an expert-supply service backed by target-derived W4 pages
and corrections. This placement is what lets the same design survive changes in GPU
generation, memory ceiling, model family, and workload. It also makes the headline
result more than a collection of isolated speedups: the measurements exercise one
portable contract for representing, supplying, and authorizing MoE state.
